% release notes for GRIDGEN version 1.1.00
\item Version 1.1.00 11/18/2026
\begin{itemize}[label=--]
  \item[] \textbf{NEW FUNCTIONALITY}
  \item A GRID\_TO\_USGDATA block also writes the MODFLOW-USG DISU array IVC and the MODFLOW 6 DISU arrays IHC, HWVA, and ANGLDEGX, so that both programs can read the connection type, face width, and face angle directly instead of computing them from the connection direction and area. ANGLDEGX includes the rotation of the grid. They are described in the Tips and Recommendations section.
  \item A new GRID\_TO\_DISVDATA block writes the input of the MODFLOW 6 DISV Package for a MODFLOW grid or a quadtree grid with the same cells in every layer: the dimensions, vertices, and CELL2D records, the TOP, BOTM, and IDOMAIN arrays, and the horizontal ghost nodes by layer and cell. Unlike FloPy's get\_gridprops\_disv, it writes grids with an active domain in some layers, and it can write inactive cells as vertical pass-through cells. It is described in the Tips and Recommendations section.
  \item A new GRID\_TO\_GSF block writes a MODFLOW-USG grid specification file, and optionally a table of the base grid cell that contains each node, from the vertices that neighboring cells share, at full precision, with or without the corners of smaller neighbors on cell edges; the GRIDGEN2GSF program of the PEST Groundwater Utilities is no longer needed. A GRID\_TO\_USGDATA block now also writes the vertices of the grid and the CELL2D records of each cell for MODFLOW 6 DISU input; DISV input is written by a GRID\_TO\_DISVDATA block. Both are described in the Tips and Recommendations section.
  \item A GRID\_TO\_USGDATA block with the new record GNC\_VERTICAL = True also writes ghost nodes for the vertical connections between a cell and the smaller cells above or below it, which correct the flow between layers with different refinement. The ghost nodes are added to the .gnc.dat file after the horizontal ones, in the same format, and are described in the Tips and Recommendations section. The default output is unchanged.
  \item[] \textbf{BUG FIXES AND OTHER CHANGES TO EXISTING FUNCTIONALITY}
  \item GRIDGEN now stops with an error message and an exit code of 1 when the definition file is missing a required record, such as NROW, names an array file that cannot be opened or an ESRI ASCII grid that cannot be read, or when ghost nodes cannot be computed. Previously, a missing record or an unreadable ESRI ASCII grid stopped the program with an assertion failure, an array file that could not be opened did not stop the run, and a ghost node error ended the run with an exit code of 0, as if it had succeeded. Output for valid input is unchanged.
  \item Calculations are compiled to be evaluated as written, without the rearrangements that some compilers made by default, such as combining a multiplication and an addition in one step. Some values in the output change in the last digit with the Intel compilers, and with gcc and clang on computers whose processors can combine the two, such as Apple silicon.
  \item Some records were read but had no effect, without a message: OUTPUT\_VTKFILE in a GRID\_INTERSECTION block, and NODE\_NUMBERING and WITHOUT\_INACTIVE in GRID\_TO\_SHAPEFILE and GRID\_TO\_VTKFILE blocks for a QUADTREE grid, which keeps the node numbering set in its QUADTREE\_BUILDER block. GRIDGEN now warns that these records are ignored. The output is unchanged.
  \item GRIDGEN now warns when a shapefile feature named in a REFINEMENT\_FEATURES, ACTIVE\_DOMAIN, or GRID\_INTERSECTION block does not intersect the grid, when a shapefile has no features of its FEATURE\_TYPE, and when a GRID\_INTERSECTION block writes no output because nothing intersects the grid. Previously, these features were ignored without a message. GRIDGEN now stops with an error when a shapefile cannot be opened, when the field named by REFINEMENT\_LEVEL\_BY\_ATTRIBUTE is not in the shapefile, when FEATURE\_TYPE is not line, polygon, or point, when an OPEN/CLOSE array file has fewer numbers than the array, when a record has no value, when a block ends with a different type than it began, and when a definition file or output file cannot be opened. Previously, most of these printed a message and continued without the input, or left array values undefined. FEATURE\_TYPE is no longer case sensitive in ACTIVE\_DOMAIN and GRID\_INTERSECTION blocks, and a shapefile with no attribute fields is now read; previously, its features were ignored.
  \item Error and warning messages now name the block, record, or file they concern in the terms of the definition file, rather than the program function that reports them; for example, a missing record is reported as ``MODFLOW\_GRID block mfg has no NROW record'', and a block that names an undefined block, or one of the wrong type, is reported with the record that names it. GRIDGEN now stops with an error when a QUADTREE structure file has a quadrant code other than 1 to 4, which was skipped without a message, when the first part of a shapefile shape does not start at its first vertex, when a feature has a coordinate that is not a number, and when GRID\_TO\_SHAPEFILE cannot create the shapefile for a MODFLOW grid or for lines; previously, it printed a message, or nothing, and continued. Output for valid input is unchanged.
  \item Multipatch shapefiles were read incorrectly, so an active domain or refinement feature from a multipatch shapefile could lose its polygons or its holes; an active domain of outer and inner rings made the whole layer inactive. Outer and first rings now start a polygon and inner rings are its holes; a plain ring starts a polygon unless it follows another ring, when it is a hole. Triangle strips and fans are skipped with a warning.
  \item Running GRIDGEN with a block that is not an action, such as a MODFLOW\_GRID, QUADTREE, REFINEMENT\_FEATURES, or ACTIVE\_DOMAIN block, wrote nothing and gave no message, as if the run had succeeded. As the GRIDGEN report describes, such a run still writes no output, but GRIDGEN now warns that the block is not an action and names the block types that write output: QUADTREE\_BUILDER, GRID\_TO\_SHAPEFILE, GRID\_TO\_USGDATA, GRID\_TO\_DISVDATA, GRID\_TO\_VTKFILE, GRID\_TO\_GSF, and GRID\_INTERSECTION.
  \item A polygon shape with more than one outer ring was read as one polygon whose later rings were holes, so a REFINEMENT\_FEATURES block refined only part of it. Each ring is now an outer boundary or a hole according to whether it lies inside another ring of the shape, whatever the order and orientation of the rings, and each hole belongs to the ring that contains it. Shapes with one outer ring and its holes are read as before.
  \item A REFINEMENT\_ARRAY (or REFINE\_BY\_ARRAY) LAYER n record in a QUADTREE\_BUILDER block, which the GRIDGEN report does not describe, was read but never used, so the grid was not refined and no message said so. GRIDGEN now warns that the record is ignored and no longer reads the array. Refine with REFINEMENT\_FEATURES instead.
  \item To find the holes of a polygon shape, each ring is now tested only against the rings whose extent contains it, so a shape with many rings, such as a coastline or a layer of water bodies dissolved into one shape, is read quickly. A shape with a million separate rings is read in under a second. The rings that are read are unchanged.
  \item With REFINEMENT\_LEVEL\_BY\_ATTRIBUTE, a line shape with several parts, a multipoint shape, a multipatch shape with several outer rings, or a null shape gave the features after it in the shapefile the refinement level of another shape, or no refinement. Each feature is now refined to the level of the shape it belongs to, so a grid refined by such a shapefile can differ from before.
  \item A line refinement feature was not refined where a segment of the line crossed the grid with both ends outside it. The cells such a segment crosses are now refined, so a grid refined by such a line can have more cells than before.
  \item A line refinement feature that ended on a cell edge or at a cell corner also refined the cells on the other side of that edge or corner, which the line does not enter. Only the cells a line crosses, passes through a corner of, or runs along an edge of are now refined. A grid refined by lines whose end points fall on cell edges or corners can have fewer cells than before.
  \item The warning about cells with negative thickness named the last such cell, counted from zero, as the one with the largest negative thickness. It now names the cell with the largest negative thickness, by its node number in the output files, with its layer, and spells thickness correctly.
  \item The GRID\_TO\_SHAPEFILE record that sets whether each cell polygon ends with its first vertex is now named CLOSE\_POLYGONS and described in the Tips and Recommendations section. Its previous, misspelled name CONCIDE is still read, with a warning that it is deprecated; when both are given, CLOSE\_POLYGONS is used.
  \item A grid definition file written to another folder than the one GRIDGEN is run from, by giving a folder in GRID\_DEFINITION\_FILE, could not be read with LOAD, and neither could one whose MODFLOW\_GRID block read its arrays from another folder. The file names in it are now relative to the folder of the grid definition file, and the QUADTREE block is named for the file without the folder, so a later GRID record names the grid as it would in the folder GRIDGEN is run from. Grid definition files written to that folder are unchanged.
  \item When GRID\_DEFINITION\_FILE named a file without an extension and its path had a period before the file name, as in out.v2/grid02qtg or ../grid02qtg, the grid was named for the part of the path before that period, and its tree structure file and top and bottom arrays were written under that name in another folder. The grid is now named for the file name without its extension, and its files are written beside the grid definition file. Names that end with an extension are unchanged.
  \item Some horizontal ghost nodes in the .gnc.dat file did not reproduce a head that varies linearly, so the flow between the two connected cells was slightly in error: those interpolated from a single cell larger than the cell that contains the ghost node and offset along the connection, and those along the edge of an active domain or of the grid, where the cells beside the cell that contains the ghost node are inactive, missing, or coarser. Such a ghost node is now interpolated from active cells nearby, with contributing factors that are not negative. In a test grid with active domains in two layers, the 1,690 of 25,598 ghost nodes that did not reproduce such a head now do, and 45 more connections have a ghost node. Where the head cannot be interpolated from the active cells nearby, the ghost node is missing or written as before, and GRIDGEN warns with the number of such connections from a cell next to an inactive cell.
  \item With VERTICAL\_PASS\_THROUGH = True, a cell above or below an inactive cell that is larger than it was also connected to the active cells on the other side of the inactive cell that are not below or above it, each over the area of the smallest cell in the column, and a connection to a cell that is below or above it could be given its flow area more than once. A cell is now connected only to the cells below or above it, over the area they share through the inactive cells.
  \item GIS software read the integer fields of the shapefiles GRIDGEN writes as double-precision fields, so another table could not be joined to them on one of these fields: the nodenumber, layer, row, and col fields of a GRID\_TO\_SHAPEFILE block, and the nodenumber field, the pointid, arcid, or polyid field, and the integer attributes copied from the input shapefile in the OUTPUT\_SHAPEFILE of a GRID\_INTERSECTION block. These fields are now 10 characters wide, which ArcGIS reads as a 32-bit integer and GDAL-based software such as QGIS and geopandas as a 64-bit integer. Their values are unchanged.
  \item NODE\_NUMBERING = zero-based in a GRID\_TO\_VTKFILE block for a MODFLOW\_GRID was ignored, so the node, layer, row, and column numbers in the VTK file always started at 1. They now start at 0 when NODE\_NUMBERING = zero-based is given. VTK files written with the default numbering are unchanged.
  \item A GRID\_INTERSECTION with FEATURE\_TYPE = polygon could list part of a hole in a cell as a separate piece of the polygon, adding its area to the cell instead of subtracting it, when an edge of the outer boundary spanned a very small angle as seen from the hole. The hole is now subtracted, and each cell is listed with the area of the polygon outside the hole.
  \item A GRID\_INTERSECTION with a polygon that has a hole counted the hole as part of the polygon: cells inside the hole were listed with their full area, and cells crossing the edge of the hole were listed twice, once with the area outside the hole and once with the area inside it. Cells inside the hole are no longer listed, and each cell is listed with the area of the polygon outside the hole. Intersections with polygons without holes are unchanged.
  \item The OUTPUT\_SHAPEFILE of a GRID\_INTERSECTION with FEATURE\_TYPE = polygon wrote the outer boundary of each polygon counterclockwise and its holes clockwise, the reverse of the shapefile specification, and did not repeat the first vertex at the end of each boundary. GIS programs warned of unclosed rings and read a cell that contains a whole hole as invalid, with a negative area. The boundaries are now written in the order the specification requires and closed, and the areas in OUTPUT\_FILE are unchanged.
  \item A GRID\_INTERSECTION with FEATURE\_TYPE = polygon could stop with an error, or report incorrect areas that changed from run to run, when vertices of the polygon or of one of its holes lay exactly on grid lines, including an edge along a grid line and vertices that lie on a grid line of a rotated grid only to within rounding. The intersection now completes with the correct areas, and gives the same result every run. A GRID\_INTERSECTION with FEATURE\_TYPE = line, or with a polygon, also stopped with an error when a line ran along a grid line and ended partway along a cell edge; it now completes. Intersections with features that do not touch grid lines are unchanged.
  \item The Linux executable now includes the C++ runtime library, so it runs on systems with an older one, such as Red Hat Enterprise Linux 9, where it stopped with a message that a version of GLIBCXX was not found. It still needs glibc 2.34 or later.
  \item[] \textbf{BUILD, TESTING, AND INTERNAL CHANGES}
  \item GRIDGEN is compiled to the C++17 standard rather than C++14, so building it from source needs a compiler that supports C++17. Its output is unchanged.
  \item The Windows executable is compiled with the LLVM-based Intel oneAPI C++ compiler, as the Linux executable already was, rather than the Microsoft compiler, and the macOS executable with gcc 15 rather than gcc 13. GRIDGEN is no longer built or tested with the classic Intel C++ compiler, which Intel has discontinued; it is built and tested with GCC, Clang, and the LLVM-based Intel compilers.
  \item The makefile in the src folder, which had not been updated since version 1.0.02 and no longer built GRIDGEN, was removed; GRIDGEN is built with meson.
  \item[] \textbf{DOCUMENTATION}
  \item Table 10 of the GRIDGEN report gives the default of SHARE\_VERTEX in a GRID\_TO\_VTKFILE block as False, but the default is True: cell vertices are shared among neighboring cells and their elevations are interpolated unless SHARE\_VERTEX is given a value other than True (in any letter case), which is read as False.
  \item A new section, Using GRIDGEN Output with MODFLOW-USG and MODFLOW 6, has a table of the file that holds each item of the MODFLOW-USG DISU, MODFLOW 6 DISU, and MODFLOW 6 DISV input, in place of a paragraph in the Introduction. The Abstract, the test problems, and the Tips and Recommendations describe the output for both programs.
\end{itemize}
